\scanpage{31}
\mspage{16}

\uline{Preuve} : (I) \quad D'après (6.3), il n'y a en effet pas plus de deux automorphismes.

(II) et (III) \quad Lorsque les courbes sont mises sous la forme (2.1), il résulte de (2.5) que les automorphismes sont décrits parleur action sur $\pi$ et que $\lambda\in G_m$ défini un automorphisme si et seulement si
\[
c_4=\lambda^4\,c_4 \qquad\text{et}\qquad c_6=\lambda^6\,c_6 \, ,
\]
d'où les assertions.

(IV) \quad Si la courbe est mise sous forme (6.1), les automorphismes sont donnés par (1.2) : $x'\handwritten{\mathrel{\tikz[baseline=-.32ex]{\draw[->,thin] (0,0) .. controls (.19em,.39ex) and (.28em,.6ex) .. (.57em,.18ex) .. controls (.78em,-.23ex) and (1.02em,-.08ex) .. (1.28em,0) -- (2.28em,.04ex);}}}x$, pour $s=t=0$ et
\[
\handwritten{\left\{}\begin{aligned}
&u^4=1\\
&r^3+ra_4+a_6(1-u^2)=0
\end{aligned}\right.
\]
et la loi de composition par
\[
(u',r')\quad (u,r)=(u\,u',u^{\raisebox{-.38ex}{$\scriptstyle\prime$}\,2}r+r')
\]

Lorsque la courbe est mise sous forme (6.7), ces équations se réduisent à $u\in\mu_4$ et $r\in F_3$ .

(V) \quad Si la courbe est mise sous forme (6.3), les automorphismes sont donnés par (1.2) : $x'\handwritten{\mathrel{\tikz[baseline=-.32ex]{\draw[->,thin] (0,0) .. controls (.14em,.37ex) and (.25em,.52ex) .. (.44em,.13ex) .. controls (.62em,-.25ex) and (.86em,-.2ex) .. (1.09em,.08ex) .. controls (1.3em,.31ex) and (1.52em,.23ex) .. (1.67em,.20ex) -- (2.36em,.16ex);}}}x$, pour $r=s^2$ et
\[
\handwritten{\left\{}\begin{aligned}
&u^3=1\\
&s^4+sa_3+a_4(1-u)=0\\
&t^2+ta_3+s^6+s^2\,a_4=0
\end{aligned}\right.
\]%
% archive sheet number: 30
